\section{Weil classes of every CM field}\label{sec:cmfields}

The Weil classes of an imaginary quadratic field are the case $m=1$ of a
construction that makes sense for every CM field $F$ of degree $2m$, and the
question whether those classes are algebraic might seem to lie outside the
methods of this paper. It does not. Every theorem of \Cref{sec:construction} that turns the conjecture for a Weil
family into a propagation statement is a theorem about a connected period
domain, a flat class on it, a group whose rational points act with dense
orbits, and a base point where the class is a polynomial in divisor classes;
none of the four uses that the field is quadratic. This section carries the
four over, and supplies the base point, which is the one input that is new.
Nothing here is conditional on a question about semiregularity: the base
point is a CM point, and the classes there are products of divisor classes
by the Lefschetz theorem on $(1,1)$ classes.%
\footnote{Nothing is claimed about the whole Hodge ring at a CM point, which
can contain classes that are not polynomials in divisor classes. What is used
is narrower: the particular Weil classes of \Cref{thm:cmbasepoint} are written
out as polynomials with rational coefficients in divisor classes, each divisor
class is algebraic by the Lefschetz theorem, and products of algebraic classes
are algebraic.}

\subsection{The families}

\begin{setupenv}[Weil type for a CM field]\label{setup:cmfield}
Let $F$ be a CM field of degree $2m$ with maximal totally real subfield
$F_{0}$, and write $\alpha\mapsto\bbar\alpha$ for the nontrivial automorphism
of $F/F_{0}$, which is complex conjugation under every embedding. The real
places of $F_{0}$ are $\tau_{1},\dots,\tau_{m}$, and each extends to a
conjugate pair $\sigma_{\tau},\bbar\sigma_{\tau}$ of embeddings of $F$ into
$\CC$. Let $V$ be an $F$-vector space of dimension $2n$, let
$H\colon V\times V\to F$ be a nondegenerate hermitian form, $F$-linear in the
first variable, and fix a totally imaginary $\xi\in F^{\times}$, that is
$\bbar\xi=-\xi$. Put
\[
  E(x,y)=\Tr_{F/\QQ}\bigl(\xi\,H(x,y)\bigr),
\]
a nondegenerate alternating $\QQ$-bilinear form with
$E(\alpha x,y)=E(x,\bbar\alpha y)$. For each $\tau$ let
$V_{\tau}=V\otimes_{F_{0},\tau}\RR$, a vector space of dimension $2n$ over
$F\otimes_{F_{0},\tau}\RR\cong\CC$, and let $(p_{\tau},q_{\tau})$ be the
signature of the hermitian form $H_{\tau}$ induced by $H$ on it. We say that
$(V,H)$ is of \emph{$(F,n)$-Weil type} if $p_{\tau}=q_{\tau}=n$ for every
$\tau$, and we call the class of $\det H$ in
$F_{0}^{\times}/\Nm_{F/F_{0}}(F^{\times})$ its discriminant $\delta$. A
polarised abelian variety $(A,\iota,E)$ of dimension $2mn$ with
$\iota\colon F\to\End^{0}(A)$ and $H^{1}(A,\QQ)\cong V$ as $F$-modules with
polarisation $E$ is of $(F,n,\delta)$-Weil type when $(V,H)$ is; as in
\Cref{lem:Kspace} we work throughout with $V=H^{1}$, a complex structure on
$V_{\RR}$ being the same datum as one on the tangent space. The Weil
classes of $A$ are
\[
  W_{F}(A)=\textstyle\bigwedge^{2n}_{F}V\;\subseteq\;\bigwedge^{2n}_{\QQ}V
  =H^{2n}(A,\QQ),
\]
the $\QQ$-subspace whose complexification is
$\bigoplus_{\sigma}\bigwedge^{2n}V_{\sigma}$, where $V_{\sigma}=V\otimes_{F,\sigma}\CC$
runs over the $2m$ eigenspaces of $F$ on $V_{\CC}$; it is a one-dimensional
$F$-vector space, of dimension $2m$ over $\QQ$. Write
\begin{quote}
$\mathbf{W}(F,n,\delta)$: every abelian variety of $(F,n,\delta)$-Weil type
has algebraic Weil classes.
\end{quote}
The period domain $\cD_{F}$ is the set of complex structures $J$ on
$V_{\RR}$ commuting with $F$ and polarised by $E$; $G_{F}=\Res_{F_{0}/\QQ}\SU(V,H)$
acts on it; $\Gamma\subset G_{F}(\QQ)$ is an arithmetic subgroup,
$\cS_{F}=\Gamma\backslash\cD_{F}$ the quotient with $p\colon\cD_{F}\to\cS_{F}$,
$\pi\colon\cA\to\cS_{F}$ the universal family, and
$\Sigma_{F}\subseteq\cD_{F}$ the set of $s$ at which $W_{F}(A_{s})$ is
algebraic. For $m=1$ this is \Cref{setup:family} and \Cref{setup:shimura},
with $\omega_{s}$ a generator of the line.
\end{setupenv}

\begin{remark}[Every such abelian variety is a member]\label{rem:cmmember}
If $A$ is any abelian variety with a CM field $F\subseteq\End^{0}(A)$ such
that $\dim_{F}H^{1}(A,\QQ)=2n$ and $\dim V_{\sigma}^{1,0}=n$ for every
embedding $\sigma$, then $A$ is of $(F,n,\delta)$-Weil type for some
$\delta$: the Rosati involution of any polarisation restricts to a positive
involution of $F$, hence to conjugation, and the argument of
\Cref{lem:hermform} writes the polarisation as $\Tr_{F/\QQ}(\xi H)$ for a
unique hermitian $H$, whose signatures are then forced by the Hodge numbers
as in \Cref{prop:cmweil}(i). These are exactly the abelian varieties whose
Weil classes for $F$ are Hodge classes, in the sense of \cite{MZ99}, so
$\mathbf{W}(F,n,\delta)$ over all $(F,n,\delta)$ is the statement that all
Weil classes of all CM fields are algebraic.
\end{remark}

\begin{proposition}[The Weil classes of a CM field]\label{prop:cmweil}
Keep \Cref{setup:cmfield}, with $(V,H)$ of $(F,n)$-Weil type.
\begin{enumerate}[label=\textup{(\roman*)},leftmargin=2.6em,itemsep=2pt]
\item A polarised complex structure $J$ is the same as an $H_{\tau}$-orthogonal
  decomposition $V_{\tau}=V_{\tau}^{+}\oplus V_{\tau}^{-}$ for every $\tau$,
  with $H_{\tau}$ definite of the sign of $\mathrm{Im}\,\sigma_{\tau}(\xi)$ on
  $V_{\tau}^{+}$ and of the opposite sign on $V_{\tau}^{-}$; then
  $\dim V^{1,0}_{\sigma_{\tau}}=\dim V_{\tau}^{+}=n$ and
  $\dim V^{1,0}_{\bbar\sigma_{\tau}}=\dim V_{\tau}^{-}=n$, so every class in
  $W_{F}(A_{J})$ is of type $(n,n)$. $W_{F}$ is a flat subsystem of
  $R^{2n}\pi_{*}\QQ$ of rank $2m$ on which $G_{F}$ acts trivially, and
  $\cD_{F}\cong\prod_{\tau}\cD_{n,n}$ is connected of dimension $mn^{2}$.
\item If one nonzero element of $W_{F}(A)$ is algebraic, every element is.
\item If $\Hg(A_{s})=G_{F}$, the Hodge ring of $A_{s}$ is generated by
  $\NS(A_{s})_{\QQ}$ and by $W_{F}$; so the Hodge conjecture for $A_{s}$ in
  every degree is the algebraicity of $W_{F}$. For $n\ge2$ the space
  $\NS(A_{s})_{\QQ}$ is the $m$-dimensional space of classes
  $\eta_{t}=[E(t\,\cdot,\cdot)]$, $t\in F_{0}$.
\item Two members of $(F,n,\delta)$-Weil type have isomorphic $(V,H)$, so
  every one of them is a fibre of $\pi$.
\end{enumerate}
\end{proposition}

\begin{proof}
(i) A complex structure commuting with $F$ preserves each $V_{\tau}$, since
$F_{0}\otimes\RR$ acts diagonally, and is $\CC$-linear on it for the
structure $F\otimes_{F_{0},\tau}\RR\cong\CC$, so it is $\pm\sqrt{-1}$ on two
complementary $\CC$-subspaces $V_{\tau}^{\pm}$. Write
$\sigma_{\tau}(\xi)=\sqrt{-1}\,r_{\tau}$ with $r_{\tau}\in\RR^{\times}$. For
$x\in V^{\epsilon}_{\tau}$ and $y\in V^{\epsilon'}_{\tau}$ one has
$H_{\tau}(Jx,Jy)=\epsilon\epsilon'H_{\tau}(x,y)$, so
$E(Jx,Jy)=E(x,y)$ forces $E(x,y)=0$ for $\epsilon\ne\epsilon'$, and replacing
$x$ by $\sqrt{-1}\,x$ then forces $H_{\tau}(x,y)=0$: the decomposition is
orthogonal. And $E(x,Jx)=2\,\mathrm{Re}\bigl(\sqrt{-1}\,r_{\tau}\,
\overline{\epsilon\sqrt{-1}}\,H_{\tau}(x,x)\bigr)=2\epsilon r_{\tau}H_{\tau}(x,x)$,
so positivity is the sign condition stated. Conversely such a decomposition
defines a polarised $J$. Under $V_{\tau}\otimes_{\RR}\CC=V_{\sigma_{\tau}}
\oplus V_{\bbar\sigma_{\tau}}$ the space $V^{1,0}$, the $\sqrt{-1}$-eigenspace
of $J$ on $V_{\CC}$, meets $V_{\sigma_{\tau}}$ in the image of $V_{\tau}^{+}$
and $V_{\bbar\sigma_{\tau}}$ in the conjugate of the image of $V_{\tau}^{-}$,
which gives the dimensions; since $H_{\tau}$ has signature $(n,n)$ the two
spaces have dimension $n$. A generator $\alpha_{\sigma}$ of
$\bigwedge^{2n}V_{\sigma}$ is then of type $(n,n)$. The subspace
$\bigwedge^{2n}_{F}V$ of $H^{2n}(A_{s},\QQ)$ is defined by the action of $F$
on the lattice and does not depend on $s$, so it is flat; $\SU(V,H)$ acts on
$\bigwedge^{2n}_{F}V$ by the $F$-determinant, which is $1$. Finally the
choice of $V_{\tau}^{+}$ is a point of the bounded symmetric domain of
$\SU(V_{\tau},H_{\tau})\cong\SU(n,n)$, of dimension $n^{2}$, and $\cD_{F}$ is
the product over $\tau$.

(ii) For $\alpha\in\cO_{F}$ the endomorphism $\iota(\alpha)$ acts on
$\bigwedge^{2n}_{F}V$ by $\alpha^{2n}$, through an algebraic correspondence,
so the algebraic classes in $W_{F}$ form a $\QQ$-subspace stable under
multiplication by every $\alpha^{2n}$. In characteristic zero the $2n$-th
powers of the elements of a field span, over the prime field, the whole
image of the $2n$-fold multiplication map, which is the field itself; so the
$\QQ$-span of $\{\alpha^{2n}w\}$ is $Fw=W_{F}$ for every $w\ne0$.

(iii) Hodge classes are the invariants of the Hodge group. Over $\CC$ the
group is $\prod_{\tau}\SL(U_{\tau})$ with $U_{\tau}=V_{\sigma_{\tau}}$, and
$H$ identifies $V_{\bbar\sigma_{\tau}}$ with $U_{\tau}^{\vee}$, so
$\bigwedge^{*}V_{\CC}=\bigotimes_{\tau}\bigwedge^{*}(U_{\tau}\oplus U_{\tau}^{\vee})$
with the $\tau$-th factor of the group acting on the $\tau$-th tensor factor
alone. The invariants of a product group in a tensor product of
representations are the tensor product of the invariants, and the invariants
of $\SL(U)$ in $\bigwedge^{*}(U\oplus U^{\vee})$ are generated by the pairing
$\psi\in U\otimes U^{\vee}$ and the two determinants, by the computation of
\Cref{thm:hdgcount}. The classes $\eta_{t}$ are Hodge classes, being
polarisations of the abelian variety $A_{s}$ with the $F_{0}$-twisted form
$E_{t}=\Tr(\xi tH)$, their complexification is $\bigoplus_{\tau}\CC\psi_{\tau}$,
and the complexification of $W_{F}$ contains all the determinants. So the
$\QQ$-algebra generated by $\NS_{\QQ}$ and $W_{F}$ has complexification equal
to the ring of invariants, hence equals $\Hdg^{*}(A_{s})$. Divisor classes
being algebraic, the conjecture reduces to $W_{F}$, and by (ii) to one class.
For the last sentence, the invariants in $\bigwedge^{2}V_{\CC}$ are
$\bigoplus_{\tau}\CC\psi_{\tau}$ when $n\ge2$: a summand
$\bigwedge^{2}U_{\tau}$ or $\bigwedge^{2}U_{\tau}^{\vee}$ has an
$\SL(U_{\tau})$-invariant only if $\dim U_{\tau}=2n$ equals $2$, and a summand
mixing two different $\tau$ has none.

(iv) Two hermitian forms over $F/F_{0}$ of the same rank, the same
discriminant and the same signature at every real place of $F_{0}$ are
isomorphic \cite[Ch.~10, \S6]{Sch85}. So the datum $(F,n,\delta)$ determines
$(V,H)$ up to isomorphism, and the Hodge structure of any member is, by (i), a
polarised complex structure on that $V_{\RR}$, that is a point of $\cD_{F}$.
\end{proof}

\subsection{A base point in every family}

\begin{theorem}[A base point in every family, for every CM field]
\label{thm:cmbasepoint}
Let $(V,H)$ be of $(F,n)$-Weil type with discriminant $\delta$. Then $\cD_{F}$
contains a point $s_{0}$ with the following properties.
\begin{enumerate}[label=\textup{(\roman*)},leftmargin=2.6em,itemsep=2pt]
\item $A_{s_{0}}$ is isogenous to $B_{\Phi}^{\,n}\times B_{\bbar\Phi}^{\,n}$,
  where $B_{\Phi}$ is an abelian $m$-fold with complex multiplication by $F$
  of type $\Phi=\{\sigma:\mathrm{Im}\,\sigma(\xi)>0\}$ and $B_{\bbar\Phi}$ its
  conjugate; there is an $H$-orthogonal basis $e_{1},\dots,e_{2n}$ of $V$
  with $H(e_{i},e_{i})$ totally positive for $i\le n$ and totally negative
  for $i>n$, and $Fe_{i}=H^{1}(B_{\Phi},\QQ)$ for $i\le n$,
  $Fe_{i}=H^{1}(B_{\bbar\Phi},\QQ)$ for $i>n$.
\item For $i\le n$ put $i'=i+n$ and let $D_{ii'}\subseteq H^{2}(A_{s_{0}},\QQ)$
  be the $F$-balanced part of $Fe_{i}\wedge Fe_{i'}$, spanned by the classes
  \[
    \delta_{i}(f)=\sum_{k}\,(\omega_{k}f\,e_{i})\wedge(\omega_{k}^{*}\,e_{i'}),
    \qquad f\in F,
  \]
  for dual $\QQ$-bases $(\omega_{k})$, $(\omega_{k}^{*})$ of $F$ with respect
  to the trace form. Every $\delta_{i}(f)$ is a Hodge class of type $(1,1)$,
  hence a divisor class, and $D_{ii'}$ is a one-dimensional $F$-vector space
  with $\alpha\cdot\delta_{i}(f)=\delta_{i}(\alpha f)$.
\item The Weil classes are the products
  \[
    w(f)=\sum_{k_{1},\dots,k_{n-1}}
      \delta_{1}(\omega_{k_{1}}f)\cdot\delta_{2}(\omega_{k_{1}}^{*}\omega_{k_{2}})
      \cdots\delta_{n-1}(\omega_{k_{n-2}}^{*}\omega_{k_{n-1}})\cdot
      \delta_{n}(\omega_{k_{n-1}}^{*}),\qquad f\in F,
  \]
  which lie in $W_{F}(A_{s_{0}})$, span it as $f$ runs over $F$, and are
  polynomials with rational coefficients in divisor classes. In particular
  $W_{F}(A_{s_{0}})$ is algebraic, and $\Sigma_{F}\ne\emptyset$ for every $F$,
  every $n\ge1$ and every $\delta$.
\end{enumerate}
\end{theorem}

\begin{proof}
(i) Since $H_{\tau}$ has signature $(n,n)$ at every $\tau$, the discriminant
$\det H$ has sign $(-1)^{n}$ at every real place, so $b=(-1)^{n}\det H$ is
totally positive, and the diagonal form $H'=\mathrm{diag}(b,1,\dots,1,-1,\dots,-1)$
with $n$ positive and $n$ negative entries has the same dimension, the same
discriminant and the same signatures as $H$. Hermitian forms over $F/F_{0}$
are classified by these invariants \cite[Ch.~10, \S6]{Sch85}, so
$(V,H)\cong(F^{2n},H')$ and the standard basis is the required orthogonal
basis $e_{i}$, with $a_{i}=H(e_{i},e_{i})$. Write $\sigma_{\tau}(\xi)=\sqrt{-1}\,r_{\tau}$
and define $J$ on $V_{\tau}=\bigoplus_{i}\CC e_{i}$ to be
$\epsilon_{i,\tau}\sqrt{-1}$ on $\CC e_{i}$, with
$\epsilon_{i,\tau}=\mathrm{sign}(r_{\tau})\,\mathrm{sign}(\tau(a_{i}))$. The
decomposition $V_{\tau}^{+}=\bigoplus_{\epsilon_{i,\tau}=+}\CC e_{i}$,
$V_{\tau}^{-}=\bigoplus_{\epsilon_{i,\tau}=-}\CC e_{i}$ is $H_{\tau}$-orthogonal
with $H_{\tau}$ of sign $\mathrm{sign}(r_{\tau})$ on $V_{\tau}^{+}$, so
\Cref{prop:cmweil}(i) makes $J$ a point $s_{0}$ of $\cD_{F}$. On
$Fe_{i}\otimes\RR=\prod_{\tau}\CC e_{i}$ the structure $J$ is multiplication by
the element $(\epsilon_{i,\tau}\sqrt{-1})_{\tau}$ of $F\otimes_{\QQ}\RR$, so
each $Fe_{i}$ is a rational sub-Hodge structure of weight one with
$F$-action, that is the $H^{1}$ of an abelian $m$-fold with complex
multiplication by $F$, of CM type
$\Phi_{i}=\{\sigma:(Fe_{i})_{\sigma}\subseteq V^{1,0}\}$, which by the
computation in \Cref{prop:cmweil}(i) is
$\{\sigma_{\tau}^{\epsilon_{i,\tau}}\}$, where $\sigma^{+}=\sigma$ and
$\sigma^{-}=\bbar\sigma$. For $i\le n$ the $a_{i}$ are totally positive, so
$\epsilon_{i,\tau}=\mathrm{sign}(r_{\tau})$ and
$\Phi_{i}=\{\sigma:\mathrm{Im}\,\sigma(\xi)>0\}=\Phi$; for $i>n$ the signs are
reversed and $\Phi_{i}=\bbar\Phi$. So $V=\bigoplus_{i}Fe_{i}$ splits
$A_{s_{0}}$ up to isogeny as stated.

(ii) Write $L_{i,\sigma}=(Fe_{i})_{\sigma}$, a line, so that
$V_{\CC}=\bigoplus_{i,\sigma}L_{i,\sigma}$ and $L_{i,\sigma}$ has type $(1,0)$
exactly when $\sigma\in\Phi_{i}$. For dual bases,
$\sum_{k}\sigma(\omega_{k})\sigma'(\omega_{k}^{*})=1$ if $\sigma=\sigma'$ and
$0$ otherwise, because the matrices $(\sigma(\omega_{k}))$ and
$(\sigma(\omega_{k}^{*}))$ are inverse transposes by the definition of the
trace form. Hence in $Fe_{i}\otimes Fe_{i'}\otimes\CC
=\bigoplus_{\sigma,\sigma'}L_{i,\sigma}\otimes L_{i',\sigma'}$ the class
$\delta_{i}(f)$ has component $\sigma(f)\,\ell_{i,\sigma}\wedge\ell_{i',\sigma}$
on the diagonal $\sigma=\sigma'$ and zero off it, for suitable generators
$\ell$; it is $F$-balanced in the sense that $(\alpha e_{i})\wedge e_{i'}$ and
$e_{i}\wedge(\alpha e_{i'})$ agree on it, which is the identity
$\alpha\cdot\delta_{i}(f)=\delta_{i}(\alpha f)$. Since $\Phi_{i'}=\bbar\Phi_{i}$,
exactly one of $L_{i,\sigma}$, $L_{i',\sigma}$ has type $(1,0)$, so every
component, hence $\delta_{i}(f)$, has type $(1,1)$; and $\delta_{i}(f)$ is
rational, being a sum of wedges of rational vectors. By the Lefschetz theorem
on $(1,1)$ classes it is a divisor class. As $f$ varies the components
$\sigma(f)$ vary independently, so $D_{ii'}\otimes\CC=\bigoplus_{\sigma}
L_{i,\sigma}\otimes L_{i',\sigma}$ and $D_{ii'}\cong F$.

(iii) The complexification of $W_{F}$ is spanned by
$\alpha_{\sigma}=\pm\prod_{i\le n}\ell_{i,\sigma}\wedge\ell_{i',\sigma}$, a
generator of $\bigwedge^{2n}V_{\sigma}=\bigwedge^{2n}\bigl(\bigoplus_{i}L_{i,\sigma}\bigr)$.
In the product $w(f)$, the component indexed by $(\sigma_{1},\dots,\sigma_{n})$,
with $\sigma_{i}$ the embedding chosen in the $i$-th factor, carries the
coefficient
$\sigma_{1}(f)\prod_{j=1}^{n-1}\sum_{k_{j}}\sigma_{j}(\omega_{k_{j}})\sigma_{j+1}(\omega_{k_{j}}^{*})$,
which vanishes unless $\sigma_{1}=\dots=\sigma_{n}$ and then equals
$\sigma_{1}(f)$. So $w(f)=\sum_{\sigma}\sigma(f)\,\alpha_{\sigma}$ up to the
sign of the reordering, an element of $W_{F}\otimes\CC$ which is rational
because it is a rational polynomial in rational classes; as $f$ runs over
$F$ the vectors $(\sigma(f))_{\sigma}$ span $\CC^{2m}$, so the $w(f)$ span
$W_{F}$. Each $w(f)$ is a rational combination of $n$-fold products of the
divisor classes $\delta_{i}(\cdot)$, and a product of divisor classes is
algebraic. Item (XXXVI) of \Cref{sec:verification} carries out every step of (i), (ii)
and (iii) in exact arithmetic for six pairs $(F,n)$, with $F$ one of
$\QQ(\zeta_{5})$, $\QQ(\zeta_{7})$, $\QQ(\zeta_{8})$, $\QQ(\zeta_{9})$, so
$m=2$ and $m=3$, and $n=1,2,3$; for $F=\QQ(\zeta_{5})$, a cyclic quartic CM
field with no imaginary quadratic subfield, at $n=2$ it also computes the
space of rational $(1,1)$ classes, of dimension $32$, and records that the
balanced sum is necessary: the plain product $\delta_{1}(1)\delta_{2}(1)$ is
not a Weil class.
\end{proof}

\begin{remark}[The base point for a quadratic field]\label{rem:cmquadratic}
For $m=1$ the point $s_{0}$ is a member isogenous to $E^{2n}$ for an elliptic
curve $E$ with complex multiplication by $K$, at which the whole Hodge ring
is generated by divisor classes, since the Hodge group is a one-dimensional
torus. It is a second base point in every family of \Cref{sec:construction},
beside the quaternionic one of \Cref{thm:everyfamily}, and it lies in the
product locus of \Cref{thm:everyn}. The formula in (iii) is then the explicit
form of the class at that point.
\end{remark}

\subsection{The reduction and the propagation}

\begin{theorem}[The reduction and the propagation, for every CM field]
\label{thm:cmpropagation}
Fix $(F,n,\delta)$ with $n\ge2$ and keep \Cref{setup:cmfield}, and let
$\omega_{s}$ denote a generator of $W_{F}(A_{s})$ chosen flatly.
\begin{enumerate}[label=\textup{(\roman*)},leftmargin=2.6em,itemsep=2pt]
\item $\mathbf{W}(F,n,\delta)$ holds if and only if $\omega_{s}$ is algebraic
  for every $s\in\cD_{F}$, and then the Hodge conjecture holds in every
  degree for every member of the family with $\Hg(A_{s})=G_{F}$.
\item $\Sigma_{F}$ is stable under $G_{F}(\QQ)$, every $G_{F}(\QQ)$-orbit in
  $\cD_{F}$ is dense, and $\Sigma_{F}$ is dense.
\item $\Sigma_{F}=p^{-1}\bigl(\bigcup_{\underline{P},\underline{q}}
  \Sigma_{\underline{P},\underline{q}}\bigr)$ with each
  $\Sigma_{\underline{P},\underline{q}}$ Zariski closed in $\cS_{F}$, the
  union countable; either $\Sigma_{F}=\cD_{F}$, or $\Sigma_{F}$ is meagre and
  the conjecture fails on a dense $G_{\delta}$.
\item If the Hilbert data of the cycle of \Cref{thm:cmbasepoint}, transported
  along the orbit $G_{F}(\QQ)\cdot s_{0}$, takes finitely many values, then
  $\Sigma_{F}=\cD_{F}$.
\item If some perfect complex $E$ on $A_{s_{0}}$ with $\ch(E)=N\omega_{s_{0}}$,
  $N\ne0$, concentrated in degree $n$, has injective semiregularity map, then
  $\Sigma_{F}=\cD_{F}$; such complexes exist.
\end{enumerate}
\end{theorem}

\begin{proof}
(i) is \Cref{prop:cmweil}(ii) and (iii): each member has its Hodge structure
at some point of $\cD_{F}$, by (i) of that proposition, so the statement
$\mathbf{W}(F,n,\delta)$ is the algebraicity of $W_{F}$ at every point, which
by (ii) is that of one generator, and by (iii) it gives the conjecture at the
members with full Hodge group.

(ii) The stability is the argument of \Cref{lem:Ginvariance}: $g\in G_{F}(\QQ)$
carries $J$ to $gJg^{-1}$ and induces, after clearing denominators, an isogeny
$A_{J}\to A_{gJ}$ whose graph transports algebraic classes to algebraic
classes, and $g$ fixes $W_{F}$ pointwise by \Cref{prop:cmweil}(i). For the
density we show that the closure of $G_{F}(\QQ)$ in $G_{F}(\RR)
=\prod_{\tau}\SU(V_{\tau},H_{\tau})$ contains the identity component, which
acts transitively on $\cD_{F}$. The form $H$ is isotropic over $F_{0}$: it is
indefinite at every real place, its signature being $(n,n)$; at every finite
place $v$ its trace form is a quadratic form in $4n\ge8$ variables over the
local field $F_{0,v}$, hence isotropic, and a hermitian form is isotropic as
soon as its trace form is \cite[Ch.~10, \S1]{Sch85}; and hermitian forms over
the quadratic extension $F/F_{0}$ satisfy the Hasse principle
\cite[Ch.~10, \S6]{Sch85}. Let
$X=\{v\in V:H(v,v)=0\}$, the cone of isotropic vectors, a quadric cone in
the $4n$-dimensional $F_{0}$-space $V$ with a nonzero $F_{0}$-point; the
projective quadric of its lines is $F_{0}$-rational, by projection from that
point, so it is $F_{0}$-birational to a projective space and therefore
satisfies weak approximation at the archimedean places, $F_{0}$ being dense
in $F_{0}\otimes\RR=\prod_{\tau}\RR$. Hence the $F_{0}$-points of $X$ are
dense in its points over $F_{0}\otimes\RR$, which are the products
$\prod_{\tau}X_{\tau}$ of the isotropic cones of the $H_{\tau}$. For an isotropic vector $v\in V$ let $U_{v}\subseteq\SU(V,H)$ be
the unipotent radical of the stabiliser of the line $Fv$, a unipotent
$F_{0}$-group; the exponential identifies $U_{v}(F_{0})$ with an
$F_{0}$-vector space and $U_{v}(F_{0}\otimes\RR)=\prod_{\tau}U_{v,\tau}(\RR)$
with the corresponding real space, so $U_{v}(\QQ)=U_{v}(F_{0})$ is dense in
$\prod_{\tau}U_{v,\tau}(\RR)$. Hence the closure $L$ of $G_{F}(\QQ)$ contains
$\prod_{\tau}U_{v,\tau}(\RR)$ for every rational isotropic $v$, in particular
the factor $U_{v,\tau}(\RR)$ for each $\tau$ separately, and its Lie algebra
$\mathfrak{l}$ contains $\mathfrak{u}_{v,\tau}$ for every $\tau$ and every
rational isotropic $v$. The map $v\mapsto\mathfrak{u}_{v,\tau}$ is continuous
on the cone $X_{\tau}\setminus0$ and rational $v$ are dense in
$\prod_{\tau}X_{\tau}$, so the closed subspace
$\mathfrak{l}\cap\mathfrak{su}(V_{\tau},H_{\tau})$ contains
$\mathfrak{u}_{v,\tau}$ for every isotropic $v\in V_{\tau}$. The span of these
is stable under the adjoint action of $\SU(V_{\tau},H_{\tau})$, which permutes
the isotropic lines, so it is a nonzero ideal of the simple Lie algebra
$\mathfrak{su}(n,n)$, hence everything. Thus
$\mathfrak{l}\supseteq\bigoplus_{\tau}\mathfrak{su}(V_{\tau},H_{\tau})$, and
$L\supseteq G_{F}(\RR)^{+}$. Since $\Sigma_{F}\ne\emptyset$ by
\Cref{thm:cmbasepoint} and is $G_{F}(\QQ)$-stable, it contains a dense orbit.

(iii) The proofs of \Cref{thm:closedlocus} and \Cref{cor:allornothing} use
only that $\cS_{F}$ is a quasi-projective variety, which is Baily and Borel
\cite{BB66}, that the relative Hilbert schemes of $\cA/\cS_{F}$ are
projective over it \cite{Gro61}, and that $W_{F}$ is a flat subsystem
consisting of Hodge classes, which is \Cref{prop:cmweil}(i); they apply
verbatim.

(iv) is the proof of \Cref{thm:degreebound}: the orbit is dense by (ii), a
finite union of the closed sets of (iii) covering it forces one of them to be
dense and closed, hence everything; then (i).

(v) The existence is the proof of \Cref{prop:ktheory}, using the algebraic
cycle of \Cref{thm:cmbasepoint} and the isomorphism
$K_{0}(A)_{\QQ}\cong\CH^{*}(A)_{\QQ}$ \cite{Ful98}; the deformation argument
is the proof of \Cref{thm:semiregclosure}, whose only inputs are the flatness
of $\ch(E)$ as a section of Hodge classes over $\cD_{F}$, which is (i) of
\Cref{prop:cmweil}, the theorem of Buchweitz and Flenner, and (iii).
\end{proof}

\begin{remark}[What is and is not new here]\label{rem:cmwhatisnew}
The propagation statements (iv) and (v) are exactly the two forms (P1) and
(P2) of \Cref{rem:tworemain}, with $F$ in place of $K$. What this section
proves is that they are the whole of the problem for every CM field, and
that their hypothesis is never vacuous: the base point exists, its cycle is
written down, and the orbit through it is dense. The theorem of
\Cref{thm:p2forces} on what (v) forces transfers as well, with the
annihilator of a class in $W_{F}$ computed over the $2m$ eigenspaces in place
of two, and we do not repeat it. Nothing above depends on Question 11.4 of
\cite{Mar25b} or on any secant construction; the route through Weil
restriction and CM points is independent of the route through secant sheaves,
and reaches every field at once.
\end{remark}

\begin{proposition}[Composite fields]\label{prop:composite}
Let $F=KF_{0}$ with $K$ imaginary quadratic, and let $A$ be of
$(F,n,\delta)$-Weil type. Then $A$ is of $(K,-1,mn)$-Weil type for the
restricted action, its Weil classes $W_{K}(A)\subseteq H^{2mn}(A,\QQ)$ lie in
the $\QQ$-span of the $m$-fold products of elements of $W_{F}(A)$, and
$W_{K}(A)$ is algebraic whenever $W_{F}(A)$ is. In particular
$\mathbf{W}(F,n,\delta)$ implies the algebraicity of the $K$-Weil classes on
every member of the $F$-family, a closed subvariety of dimension $mn^{2}$ of
the $K$-family of dimension $m^{2}n^{2}$; the converse implication is not
available, and \Cref{thm:cmbasepoint} is not a consequence of
\Cref{thm:everyfamily}.
\end{proposition}

\begin{proof}
The eigenspace $V_{\tau_{K}}$ of $K$ on $V_{\CC}$ for the embedding
$\tau_{K}$ is $\bigoplus_{\sigma|\tau_{K}}V_{\sigma}$, of dimension $2mn$
with $\dim V_{\tau_{K}}^{1,0}=mn$ by \Cref{prop:cmweil}(i), so $A$ is of
$K$-Weil type with balanced signature, and
$\bigwedge^{2mn}V_{\tau_{K}}=\bigotimes_{\sigma|\tau_{K}}\bigwedge^{2n}V_{\sigma}$
is spanned by $\prod_{\sigma|\tau_{K}}\alpha_{\sigma}$. So
$W_{K}\otimes\CC$ is contained in the complexification of the $\QQ$-span $P$ of
the $m$-fold products of elements of $W_{F}$, and a rational subspace whose
complexification lies in that of $P$ lies in $P$. Products of algebraic
classes are algebraic. Item (XXXVI) checks the identity for $F=\QQ(\zeta_{8})$
and $K=\QQ(\sqrt{-1})$. The final sentence is dimension count: for $m\ge2$ the
$F$-family is a proper subvariety, and the $K$-Weil classes on it are products
that carry no information about their factors.
\end{proof}
